\chapter{Fix 36: Lieb-Thirring Stability for Lattice Fermions}
\label{ch:fix36_lieb_thirring}

\section{The Stability Problem}
To prove that the mass gap $\Delta(M)$ does not vanish for finite $M$, we must ensure that the inclusion of fermions does not lead to an instability where the ground state energy density diverges or the system undergoes a collapse.

\section{Lattice Lieb-Thirring Inequality}
We adapt the Lieb-Thirring inequality to the lattice gauge theory context.
\begin{theorem}[Lattice Kinetic Stability]
\label{thm:lattice_lieb_thirring}
For the Wilson-Dirac operator $H_F = \bar{\psi}(D_W+M)\psi$, the partial trace of the negative eigenvalues (filling the Dirac sea) is bounded by the kinetic energy:
\begin{equation}
\text{Tr}(H_F)_- \ge - C_d L^{-d} \int_{\Lambda^*} |p|^2 dp \sim -C \cdot Vol
\end{equation}
This implies that the fermionic contribution to the vacuum energy is extensive (proportional to volume) and bounded from below.
\end{theorem}

\section{Absence of Fermionic Condensation Instability}
We use this inequality to bound the fluctuations of the fermion condensate $\langle \bar{\psi} \psi \rangle$.
\begin{lemma}
For $M > 0$, the condensate is uniformly bounded:
\begin{equation}
|\langle \bar{\psi} \psi \rangle| \le \frac{C}{M}
\end{equation}
\end{lemma}
This bound prevents the "runaway" condensation that could theoretically close the gap at intermediate coupling.
